\section*{Devoir maison }

On s'intéresse au code suivant :

\begin{lstlisting}[language=Python]
# traduction, height : taille, weight : poids, hp : points de vie, speed : vitesse
pokemons = [{'name' : 'pikachu', 'height' : 4, "weight" : 60,
             'types' : ('electric',),
             'stats' : {'hp': 35, 'attack': 55, 'defense': 40,
                        'special-attack': 50,  'special-defense': 50,
                        'speed': 90}},
            {'name': 'dracolosse', 'height': 22, 'weight': 2100,
             'types': ('dragon', 'flying'),
             'stats': {'hp': 91, 'attack': 134, 'defense': 95,
                       'special-attack': 100, 'special-defense': 100,
                       'speed': 80}},
            {'name': 'ronflex', 'height': 21, 'weight': 4600,
             'types': ('normal',),
             'stats': {'hp': 160, 'attack': 110, 'defense': 65,
                       'special-attack': 65, 'special-defense': 110,
                       'speed': 30}}]
                       
\end{lstlisting}

\begin{Exercise}[title={Structure de données imbriquées}]
    \Question Quel est le type de la variable pokemons ?
    \Question Quel est le type de pokemons[0] ?
    \Question Quel est le type de pokemons[1]["types"] ?
    \Question Quel est le type de pokemons[1]["stats"] ?
    \Question Donner une instruction permettant de récupérer l'attaque spéciale de pikachu.\\ \emptybox{1cm}
    \Question Donner une instruction permettant de récupérer le deuxième type de dracolosse.\\ \emptybox{1cm}
    \Question Donner une instruction permettant de modifier l'attaque de ronflex pour la mettre à 1000.\\ \emptybox{1cm}
    \Question Compléter la fonction ci-dessous qui cherche le Pokémon le plus lourd dans la liste et renvoie son nom.
    \begin{lstlisting}[language=Python]
def plus_lourd(pokemons):
    #initialisation
    meilleur = ""
    meilleur_poids = 0
    for pk in pokemons:
        if ... :
            meilleur =
            meilleur_poids =
    return ...
    \end{lstlisting}
\end{Exercise}